% Generated draft: requires mathematical and source review.
\documentclass[11pt]{article}
\usepackage[margin=25mm]{geometry}
\usepackage{fontspec}
\setmainfont{Latin Modern Roman}
\usepackage{amsmath,amssymb,mathtools,unicode-math}
\setmathfont{Latin Modern Math}
\usepackage{enumitem,needspace,xurl,hyperref}
\setlength{\parindent}{0pt}
\setlength{\parskip}{5pt}
\setlength{\emergencystretch}{3em}
\newcommand{\N}{\mathbb{N}}
\newcommand{\Z}{\mathbb{Z}}
\newcommand{\Q}{\mathbb{Q}}
\newcommand{\R}{\mathbb{R}}
\newcommand{\C}{\mathbb{C}}
\newcommand{\F}{\mathbb{F}}
\newcommand{\A}{\mathbb{A}}
\newcommand{\PP}{\mathbb{P}}
\newcommand{\E}{\mathbb{E}}
\newcommand{\eps}{\varepsilon}
\newcommand{\dd}{\,\mathrm{d}}
\begin{document}
\section*{UnsolvedMath ID: 800415}
\textbf{Source problem code:} catalogue entry
\subsection{Definitions and mathematical statement}
Let \(S\) be a finite nonempty set. Write \(\operatorname{End}_{\mathsf{Set}}(S)=\{f:S\to S\}\) and let \(\operatorname{Aut}_{\mathsf{Set}}(S)\) be its subset of bijections. For \(f:S\to S\), define
\[
\operatorname{Per}(f,S)=\{x\in S: f^{\circ m}(x)=x\text{ for some integer }m\geq 1\}.
\]
The functional directed graph \(\Gamma_f\) has vertex set \(S\) and one directed edge \(x\to f(x)\) for each \(x\in S\). Interpreting components as connected components of the underlying undirected graph, the statistics mentioned in the source include \(\#\operatorname{Per}(f,S)\), the number of components, the maximum number of vertices in a component, and the maximum length of a directed cycle. This component convention makes the terminology explicit; the supplied record does not specify a convention.

\textbf{Question 38 (Vague Motivating Question).} To what extent do polynomial maps behave like random set maps? More specifically, as \(\#S\to\infty\), which statistics of periodic points and functional graphs for polynomial maps exhibit behavior similar to the corresponding statistics for random elements of \(\operatorname{End}_{\mathsf{Set}}(S)\) or \(\operatorname{Aut}_{\mathsf{Set}}(S)\)?

\textbf{Unresolved formulation qualifications.} The source presents a motivating question rather than a fully quantified conjecture. It does not specify the sets on which the polynomial maps act, a coefficient ring or field, the dimension, degree restrictions, a family or probability distribution of polynomial maps, or the limiting regime beyond \(\#S\to\infty\). It also does not specify a probability distribution on the comparison maps, which of the two comparison classes is intended in a given setting, or whether similarity means agreement of expectations, limiting distributions, orders of growth, or some other property. Consequently, the supplied record does not determine a unique precise asymptotic assertion; these choices must remain unresolved.
\subsection{Short English statement}
The question asks how closely periodic points, components, and cycles of polynomial maps resemble those of random maps on large finite sets. The supplied source explicitly calls this a vague motivating question and leaves the polynomial family, sampling rules, and meaning of statistical agreement unspecified.
\subsection{Sources}
\begin{itemize}
\item Local UnsolvedMath registry, permanent ID \(800415\), source problem code \texttt{catalogue entry}, upstream ID \(800415\), upstream revision \texttt{d371ee917ac9647fbb06cd6fc529deec6f3dda82}. The registry labels the problem open; this is a supplied status label, not an independently verified mathematical or literature assessment.
\item Benedetto, DeMarco, Ingram, Jones et al., \emph{Current Trends and Open Problems in Arithmetic Dynamics} (2018), Question 38, attributed in the record as a ``Vague Motivating Question.'' Supplied source URL: \url{https://arxiv.org/abs/1806.04980}. The record reports access on 2026-10-16 and extraction by \texttt{arxiv\_latex}; these provenance details are supplied and have not been independently verified here.
\item The source cites references labeled MR1083961, MR0119227, and MR0062973 for statistics of random maps or permutations. It specifically attributes analysis of many properties of average random maps to Flajolet and Odlyzko, reference MR1083961. Full bibliographic entries for these three references are not supplied.
\end{itemize}
\end{document}
